\documentclass[12pt]{article}

\usepackage[T1]{fontenc}
\usepackage[utf8]{inputenc}
\usepackage[russian]{babel}

\usepackage[left=2cm,right=2cm,top=2cm,bottom=2cm,bindingoffset=0cm]{geometry}

% AMS packages
\usepackage{amsmath}
\usepackage{amssymb}
\usepackage{amsfonts}
\usepackage{amsthm}

% floating option for figures
\usepackage{float}
\usepackage{graphicx}

\usepackage[autostyle]{csquotes}

\usepackage{fancyhdr}
\pagestyle{fancy}
% modifying page layout using fancyhdr
\fancyhf{}
\renewcommand{\sectionmark}[1]{\markright{\thesection\ #1}}
\renewcommand{\subsectionmark}[1]{\markright{\thesubsection\ #1}}

\rhead{\fancyplain{}{\rightmark }}
\cfoot{\fancyplain{}{\thepage }}

\usepackage{array}
\usepackage{braket}

\makeatletter
\def\env@dmatrix{\hskip -\arraycolsep
  \let\@ifnextchar\new@ifnextchar
  \extrarowheight=2ex
  \array{*\c@MaxMatrixCols{>{\displaystyle}c}}}

\newenvironment{dmatrix}
  {\env@dmatrix}
  {\endarray\hskip-\arraycolsep}

\newenvironment{bdmatrix}
  {\left[\env@dmatrix}
  {\endmatrix\right]}
\makeatother

\newcommand{\mf}{\mathbf}

\newcommand{\lb}{\left(}
\newcommand{\rb}{\right)}
\newcommand{\lsq}{\left[}
\newcommand{\rsq}{\right]}
\newcommand{\lc}{\left\{}
\newcommand{\rc}{\right\}}

\DeclareMathAlphabet{\mymathbb}{U}{BOONDOX-ds}{m}{n}
\newcommand{\bbzero}{\mymathbb{0}}
\newcommand{\bbone}{\mymathbb{1}}

\def\hh{_{\rm H}}

\usepackage[toc, page]{appendix}

\usepackage{chngcntr}
\usepackage{etoolbox}

\AtBeginEnvironment{subappendices}{%
\chapter*{Приложение 2.A}
\addcontentsline{toc}{chapter}{Приложение}
\counterwithin{figure}{section}
\counterwithin{table}{section}
}

\usepackage{mathtools}

\DeclarePairedDelimiter\abs{\lvert}{\rvert}%
\DeclarePairedDelimiter\norm{\lVert}{\rVert}%

% Swap the definition of \abs* and \norm*, so that \abs
% and \norm resizes the size of the brackets, and the 
% starred version does not.
\makeatletter
\let\oldabs\abs
\def\abs{\@ifstar{\oldabs}{\oldabs*}}
%
\let\oldnorm\norm
\def\norm{\@ifstar{\oldnorm}{\oldnorm*}}
\makeatother

\newcommand{\bnp}[1]{b_n^{(#1)}}

% hyperlinks for eqref and cite
\usepackage{hyperref}

% indentation in the first paragraph of section or paragraph
\usepackage{indentfirst}

\usepackage{dsfont}
\newcommand{\boltz}{k_\text{B}}
\newcommand{\kb}{k_\text{B}}
\newcommand{\intty}{\int\limits_{-\infty}^\infty}
\newcommand{\mean}[1]{\langle #1 \rangle}
\newcommand{\bs}{\boldsymbol}
\newcommand{\mL}{\mathcal{L}}
\newcommand{\mH}{\mathcal{H}}
\newcommand{\mN}{\mathcal{N}}
\newcommand{\mHplane}{\mathcal{H}_\text{plane}}
\newcommand{\bbS}{\mathds{S}}
\newcommand{\bxi}{\boldsymbol{\xi}}
\newcommand{\rfixed}{r_\text{fixed}}
\newcommand{\Tl}{T_\text{L}}
\newcommand{\Th}{T_\text{H}}
\newcommand{\bOmega}{\boldsymbol{\Omega}}
\newcommand{\bbI}{\mathds{I}}
\newcommand{\bba}{\mathbbm{a}}
\newcommand{\bbA}{\mathds{A}}
\newcommand{\bbB}{\mathds{B}}
\newcommand{\bbG}{\mathds{G}}
\newcommand{\bbF}{\mathds{F}}
\newcommand{\bbM}{\mathds{M}}
\newcommand{\bbV}{\mathds{V}}
\newcommand{\bbW}{\mathds{W}}
\newcommand{\EOmega}{\boldsymbol{\Omega}_{\rm \textbf{e}}}
\newcommand{\mfpe}{\mf{p_e}}
\newcommand{\mfpet}{\mf{p_e^+}}
\newcommand{\mfq}{\mf{q}}
\newcommand{\mfp}{\mf{p}}
\newcommand{\mfJ}{\mf{J}}
\newcommand{\Jx}{J_\text{X}}
\newcommand{\Jy}{J_\text{Y}}
\newcommand{\Jz}{J_\text{Z}}
\newcommand{\dJx}{\dot{\Jx}}
\newcommand{\dJy}{\dot{\Jy}}
\newcommand{\dJz}{\dot{\Jz}}
\newcommand{\bUpsilon}{\boldsymbol{\Upsilon}_e}
\newcommand{\pe}{\mathbf{p}_e}

\usepackage[raggedright]{titlesec}
\usepackage{fourier, erewhon}

\usepackage{microtype}
 \SetTracking[no ligatures={f}]{encoding=*}{100}
% changing fontsize of chapter titles
 \titleformat{\chapter}[display]
 {\bfseries\Large\lsstyle\SetTracking[no ligatures = {f}]{encoding = *}{50}\filleft}
 {\MakeUppercase{\chaptername}\enspace\thechapter}
 {2ex}
 {\titlerule[1pt]\vspace{2ex}\MakeUppercase}%
\titlespacing*{\chapter}{0pt}{-60pt}{10ex}

% changing fontsize of section titles
\titleformat{\section}
 {\normalfont\fontsize{16}{15}\bfseries}{\thesection}{1em}{}

 \usepackage{tabularx}
 \usepackage{multirow}
 \usepackage{booktabs}
 \usepackage{array}
 \usepackage{siunitx}
 \sisetup{output-exponent-marker=\ensuremath{\mathrm{e}}}

% \usepackage[backend=biber, bibencoding=auto]{biblatex}
% \usepackage[backend=biber, 
%             bibencoding=auto, 
%             style=gost-numeric,
%             language=auto,
%             sorting=none]{biblatex}
% \addbibresource{biblio.bib}

% disable per-chapter figure counter
% \usepackage{chngcntr}
% \counterwithout{figure}{chapter}
% \usepackage[figurewithin=none]{caption}

\usepackage{mattens}

\usepackage[backend=biber, style=gost-numeric, language=auto, sorting=none]{biblatex}
\bibliography{biblio}
%\bibliographystyle{utf8gost705u_my.bst}
%\bibliographystyle{unsrt}

\begin{document}

\begin{titlepage}
\centering
\textbf{\large Московский государственный университет имени М.В.\,Ломоносова\\
\vspace*{0.1cm} Химический факультет\\
\vspace*{0.1cm}
\noindent\makebox[\linewidth]{\rule{\paperwidth}{0.4pt}}
\vspace*{0.1cm}
 Кафедра физической химии\\
\vspace*{0.1cm} Лаборатория строения и квантовой механики молекул \\}

\begin{center}
\includegraphics[width=0.3\textwidth]{pictures/logo.jpg}
\end{center}

\vspace{0.5cm}
{\Large \textbf{Реферат}} \\ \vspace{0.3cm} {\large на тему} \\ \vspace{0.3cm} 
{\large \enquote{Применение формализма корреляционных функций для описания взаимодействия электромагнитного излучения в ИК области с молекулярными системами}} \\ 
\vspace{0.75cm}
{\large Финенко Артема Андреевича}

\vfill
\large\textbf{Москва\\ 2019}
\end{titlepage}



\section{Взаимодействие электромагнитного излучения с молекулярными системами}
Рассмотрим систему из $N$ взаимодействующих частиц (молекул) в квантовом состоянии $\ket{j}$. Обозначим гамильтониан системы через $\hat{H}_0$. Пусть система подвергается воздействию электрического поля частоты $\omega$, которое индуцирует переходы в другие состояния $\ket{k}$ системы при условии, что частота поля близка к частотам Бора системы
%
\begin{gather}
    \omega_{jk} = (E_j - E_k) / \hbar.
\end{gather}

Будем считать, что длина волны рассматриваемого поля $\lambda$ много больше размеров молекул в системе, поэтому в локальной окрестности молекул поле можно считать однородным в пространстве.  Электрическая составляющая плоской волны может быть записана в виде суммы
%
\begin{gather}
    \mf{E}(t) = E_0 \boldsymbol{\varepsilon} \cos \omega t = \frac{E_0 \boldsymbol{\varepsilon}}{2} \lb e^{i \omega t} + e^{- i\omega t} \rb,
\end{gather}
%
где $E_0$ -- амплитуда волны, а $\boldsymbol{\varepsilon}$ -- единичный вектор, ориентированный вдоль направления распространения волны. Энергия взаимодействия молекулярной системы с электрическим полем в дипольном приближении равна
%
\begin{gather}
    W(t) = - (\boldsymbol{\mu} \cdot \mf{E}) = - \frac{E_0}{2} \lb \boldsymbol{\mu} \cdot \boldsymbol{\varepsilon} \rb \lsq e^{i \omega t} + e^{-i \omega t} \rsq, \label{part1-dipole-approximation} 
\end{gather}
%
где через $\boldsymbol{\mu}$ обозначен полный дипольный момент системы. Взаимодействие молекулярных систем с электромагнитным полем часто рассматривают в этом приближении, считая поле классическим объектом. \par  
Взаимодействие молекулярных систем с электрическим полем традиционно рассматривается в рамках временной теории возмущений \cite{cohentanuji, greiner}. Вероятность индуцированного возмущением $W(t)$ перехода между состояниями невозмущенной системы $\ket{j} \rightarrow \ket{k}$ в первом порядке временной теории возмущения равна \cite{cohentanuji}
%
\begin{gather}
    \mathcal{P}_{jk}(t) = \frac{1}{\hbar^2} \abs{ \int\limits_0^t W_{jk}(t^\prime) e^{i \omega_{jk} t^\prime} dt^\prime }^2,
\end{gather}
%
где через $W_{jk}(t)$ обозначен матричный элемент возмущения на состояниях невозмущенной системы, равный в данном случае
%
\begin{gather}
    W_{jk} = -\frac{E_0}{2} \bra{j} \boldsymbol{\mu} \cdot \boldsymbol{\varepsilon} \ket{k} \lsq e^{i \omega t} + e^{-i \omega t} \rsq.
\end{gather}

Коэффициенты разложения первого порядка $b_n^{(1)}(t)$ возмущенной волновой функции $\ket{\psi(t)}$ в базисе собственных функций невозмущенного гамильтониана равны 
\begin{gather}
    b_{n}^{(1)}(t; \omega) = -\frac{E_0}{2 \hbar} \bra{j} \boldsymbol{\mu} \cdot \boldsymbol{\varepsilon} \ket{k} \lsq \frac{e^{i (\omega_{jk} + \omega) t} - 1}{\omega_{jk} + \omega} + \frac{e^{i (\omega_{jk} - \omega) t} - 1}{\omega_{jk} - \omega} \rsq. \label{part1-bn-expression}
\end{gather}
%
Квадрат модуля коэффициента $b_{n}^{(1)}$ определяет вероятность перехода в $n$-ое стационарное состояние невозмущенного гамильтониана. Возводя в квадрат выражение \eqref{part1-bn-expression}, после алгебраических преобразований приходим к
%
\begin{gather}
    \hspace*{-0.5cm}
    \abs{b_{n}^{(1)}(t)}^2 = \frac{E_0}{\hbar^2} \lsq \frac{\sin^2 \lb \frac{1}{2} \lb \omega_{jk} + \omega \rb t \rb}{\lb \omega_{jk} + \omega \rb^2} + \frac{\sin^2 \lb \frac{1}{2} \lb \omega_{jk} - \omega \rb t \rb}{\lb \omega_{jk} - \omega \rb^2} + \frac{ 8 \cos \lb \omega t \rb \sin \lb \frac{1}{2} \lb \omega_{jk} + \omega \rb t \rb \sin \lb \frac{1}{2} \lb \omega_{jk} - \omega \rb t \rb}{\lb \omega_{jk} + \omega \rb \lb \omega_{jk} - \omega \rb} \rsq. \label{part1-bn-expression2}
\end{gather}
 
В данном контексте нас интересуют не вероятности переходов, а скорости переходов $\Gamma_{jk}$ (иными словами, вероятности $P_{jk}$, отнесенные к единице времени) при больших значениях $t$
%
\begin{gather}
    \Gamma_{jk} = \lim_{t \rightarrow \infty} \frac{P_{jk}}{t}.
\end{gather}
%
При предельном переходе в \eqref{part1-bn-expression2} первые два слагаемые преобразуются к дельта-функ\-циям, в то время как последнее слагаемое оказывается нулевым \cite{baym-quantum-mechanics}. Итак, выражение для скорости переходов оказывается следующим \cite{mcquarrie-statistical-mechanics}  
%
\begin{gather}
    \Gamma_{jk} = \frac{\pi E_0^2}{2 \hbar^2} \abs{ \bra{j} \boldsymbol{\mu} \cdot \boldsymbol{\varepsilon} \ket{k} }^2 \Big[ \delta \lb \omega_{jk} - \omega \rb + \delta \lb \omega_{jk} + \omega \rb \Big]. \label{part1-transition-rate} 
\end{gather}

Выражение \eqref{part1-transition-rate} определяет скорость переходов между конкретными состояниями невозмущенной системы $\ket{j}$ и $\ket{k}$. Скорость поглощения энергии в ходе переходов между этими состояниями равна $\hbar \omega_{jk} \Gamma_{jk}$, т.к. в процессе одного акта поглощения система поглощает энергию, равную $\hbar \omega_{jk}$. Скорость поглощения энергии в ходе переходов с заданного уровня $\ket{j}$ может быть получена при суммировании по всем состояниям $\ket{k}$, которые доступны системе для перехода. Наконец, суммарная скорость поглощения энергии излучения системой $-\dot{E}_\text{rad}$ получается в результате суммирования по всем возможным начальным состояниям $\ket{j}$ с соответствующими заселенностями $\rho_j$
\begin{gather}
    -\dot{E}_\text{rad} = \sum_j \sum_k \rho_j \hbar \omega_{jk} \Gamma_{jk} = \frac{\pi E_0^2}{2 \hbar} \sum_{j, k} \omega_{jk} \rho_j \abs{ \bra{j} \boldsymbol{\mu} \cdot \boldsymbol{\varepsilon} \ket{k} }^2 \Big[ \delta \lb \omega_{jk} - \omega \rb + \delta \lb \omega_{jk} + \omega \rb \Big]. \label{part1-energy-absorption-rate}
\end{gather}

Для получения более симметричной формы уравнения \eqref{part1-energy-absorption-rate} осуществим алгебраические преобразования. Рассмотрим отдельно вторую сумму, получающуюся при раскрытии скобок в уравнении \eqref{part1-energy-absorption-rate}. Поменяем местами индексы $j \leftrightarrow k$, что обосновывается тем, что оба индекса пробегают по всем квантовым состояниям системы,
\begin{gather}
    \frac{\pi E_0^2}{2 \hbar} \sum_{j, k} \omega_{jk} \rho_j \abs{ \bra{j} \boldsymbol{\mu} \cdot \boldsymbol{\varepsilon} \ket{k} }^2 \delta \lb \omega_{jk} + \omega \rb = -\frac{\pi E_0^2}{2 \hbar} \sum_{j, k} \omega_{jk} \rho_j \abs{ \bra{j} \boldsymbol{\mu} \cdot \boldsymbol{\varepsilon} \ket{k} }^2 \delta \lb \omega_{jk} - \omega \rb. \label{part1-index-change}
\end{gather}

Подстановка \eqref{part1-index-change} в \eqref{part1-energy-absorption-rate} приводит к выражению, в котором индексы $j$ и $k$ входят симметричным образом 
%
\begin{gather}
    -\dot{E}_\text{rad} = \frac{\pi E_0^2}{2 \hbar} \sum_{j, k} \omega_{j k} \lb \rho_j - \rho_k \rb \abs{ \bra{j} \boldsymbol{\mu} \cdot \boldsymbol{\varepsilon} \ket{k} }^2 \delta \lb \omega_{jk} - \omega \rb.
\end{gather}

Т.к. мы предполагаем, что возмущение достаточно слабо и действует на протяжении малого промежутка времени, то будем считать, что в любой момент система находится в состоянии теплового равновесия при температуре $T$. Используя это предположение, выразим заселенность $k$-ого состояния через заселенность $j$-го состояния (понятно, что можно воспользоваться и обратной связью, т.к. мы специально привели формулу к симметричному относительно замены индексов виду)
%
\begin{gather}
    \rho_k = \rho_j \exp \lb - \beta \hbar \omega_{jk} \rb,
\end{gather}
%
где $\beta = \boltz T$ и $\boltz$ -- постоянная Больцмана. Кроме того, вследствие того, что внутри суммы находятся дельта-функции, центрированные на $\omega_{jk}$ , функции от $\omega$, вычисленные при частоте $\omega_{jk}$, могут быть вынесены из под знака суммы: 
\begin{align}
    -\dot{E}_\text{rad} &= \frac{\pi E_0^2}{2 \hbar} \sum_{j, k} \omega_{jk} \rho_j \lb 1 - \exp \lb - \beta \hbar \omega_{jk} \rb \rb \abs{ \bra{j} \boldsymbol{\mu} \cdot \boldsymbol{\varepsilon} \ket{k} }^2 \delta \lb \omega_{jk} - \omega \rb = \\
    &= \frac{\pi E_0^2}{2 \hbar} \omega \lb 1 - \exp \lb - \beta \hbar \omega \rb \rb \sum_{j, k} \rho_j \abs{ \bra{j} \boldsymbol{\mu} \cdot \boldsymbol{\varepsilon} \ket{k} }^2 \delta \lb \omega_{jk} - \omega \rb. 
\end{align}

Суммарный поток энергии $I$, переносимой электромагнитной волной через среду с показателем преломления $n$, равен усредненному по времени модулю вектора Пойнтинга $\langle S \rangle$ и равен \cite{mcquarrie-statistical-mechanics}
\begin{gather}
    I = \langle S \rangle = \frac{c}{8 \pi} n E_0^2,
\end{gather}
%
где $c$ -- скорость света в вакууме. Показатель поглощения среды $\alpha(\omega)$ определяют как отношение энергии, поглощаемой средой в единицу времени при частоте $\omega$, к энергии, переносимой электромагнитной волной в единицу времени \cite{mcquarrie-statistical-mechanics}
\begin{gather}
    \alpha(\omega) = \frac{-\dot{E}_\text{rad}}{I} = \frac{4 \pi^2}{\hbar c n} \omega \lb 1 - e^{- \beta \hbar \omega} \rb \sum_{j, k} \rho_j \abs{ \bra{j} \boldsymbol{\mu} \cdot \boldsymbol{\varepsilon} \ket{k}}^2 \delta \lb \omega_{jk} - \omega \rb. \label{part1-absorption-coefficient-definition}
\end{gather}

На основании выражения \eqref{part1-absorption-coefficient-definition} определяют спектральную функцию $J(\omega)$ \cite{gordon1968}
\begin{gather}
    J(\omega) = \frac{3 \hbar c n \alpha(\omega)}{4 \pi^2 \omega \lb 1 - e^{-\beta \hbar \omega} \rb} = 3 \sum_{j, k} \rho_j \abs{ \bra{j} \boldsymbol{\mu} \cdot \boldsymbol{\varepsilon} \ket{k} }^2 \delta \lb \omega_{jk} - \omega \rb. \label{part1-spectral-function-definition}
\end{gather}

Отметим, что обычно спектральную функцию определяют таким образом, чтобы ее интеграл по частотному диапазону был равен единице. Здесь принято несколько иное определение, эта функция не нормирована. Кроме того, спектральная функция может быть определена и для отрицательных частот, в этом случае она относится к испусканию излучения. \par
Альтернативную форму выражения \eqref{part1-spectral-function-definition} получают, осуществляя смену Шредингеровского представления квантовой механики на представление Гейзенберга. Состояния в представлении Гейзенберга не зависят от времени -- временная эволюция заложена в операторах. Эволюция оператора $\hat{A}(t)$ описывается оператором эволюции $\hat{U}(t)$
\begin{gather}
    \hat{A}(t) = \hat{U}^{+}(t) \hat{A}(0) \hat{U}(t) = e^{\frac{i}{\hbar} \hat{H} t} \hat{A}(0) e^{-\frac{i}{\hbar} \hat{H} t}. 
\end{gather}

Удобно перейти в выражении \eqref{part1-spectral-function-definition} к представлению Гейзенберга, представив дельта-функцию как Фурье-образ мнимой экспоненты
\begin{gather}
    \delta (\omega) = \frac{1}{2\pi} \int\limits_{-\infty}^\infty e^{i \omega t} dt,
\end{gather}
получаем
\begin{gather}
    J(\omega) = \frac{3}{2\pi} \sum_{j,k} \rho_j \bra{j} \boldsymbol{\mu} \cdot \bs{\varepsilon} \ket{k} \bra{k} \boldsymbol{\mu} \cdot \boldsymbol{\varepsilon} \ket{j} \int\limits_{-\infty}^\infty \exp \lsq \lb \frac{E_j - E_k}{\hbar} - \omega \rb t \rsq dt. \label{part1-spectral-function-2}
\end{gather}

Т.к. состояния $\ket{k}$ и $\ket{j}$ являются собственными состояниями гамильтониана $\hat{H}_0$, то
\begin{gather}
    \exp \lb -\frac{i}{\hbar} E_j t \rb \ket{j} = \exp \lb -\frac{i}{\hbar} \hat{H}_0 t \rb \ket{j}, \quad \exp \lb \frac{i}{\hbar} E_k t \rb \bra{k} = \exp \lb \frac{i}{\hbar} \hat{H}_0 t \rb \bra{k}. 
\end{gather}

Произведение матричного элемента и экспоненты с Боровской частотой, получаемое при внесении матричного элемента под интеграл в \eqref{part1-spectral-function-2}, легко переводится в представление Гейзенберга
\begin{gather}
    \exp \lb \frac{i}{\hbar} \lb E_k - E_j \rb t \rb \bra{k} \boldsymbol{\mu} \cdot \boldsymbol{\varepsilon} \ket{j} = \bra{k} \exp \lb \frac{i}{\hbar} \hat{H}_0 t \rb \boldsymbol{\mu} \cdot \boldsymbol{\varepsilon} \exp \lb -\frac{i}{\hbar} \hat{H}_0 t \rb \ket{j} = \bra{k} \boldsymbol{\mu}(t) \cdot \boldsymbol{\varepsilon} \ket{j}. \label{part1-matrix-element}
\end{gather}

Подставляя \eqref{part1-matrix-element} в \eqref{part1-spectral-function-2}, приходим к
\begin{gather}
    J(\omega) = \frac{3}{2\pi} \int\limits_{-\infty}^\infty \sum_{j,k} \rho_j \bra{j} \boldsymbol{\mu}(0) \cdot \boldsymbol{\varepsilon} \ket{k} \bra{k} \boldsymbol{\mu}(t) \cdot \boldsymbol{\varepsilon} \ket{j} e^{-i \omega t} dt.
\end{gather}

Учитывая соотношение замкнутости \eqref{app-completeness}, получим
\begin{gather}
    J(\omega) = \frac{3}{2\pi} \int\limits_{-\infty}^\infty \sum_j  \rho_{j} \bra{j} \boldsymbol{\mu}(0) \cdot \boldsymbol{\varepsilon} \cdot \boldsymbol{\mu}(t) \cdot \boldsymbol{\varepsilon} \ket{j} e^{-i\omega t} dt.
\end{gather}

Сумма в подынтегральном выражении является квантово-механическим средним по ансамблю значением оператора, которое в дальнейшем будем обозначать через $\langle \cdots \rangle$. Считая среду изотропной, проинтегрируем по всем возможным ориентациям $\boldsymbol{\varepsilon}$:
\begin{gather}
    J(\omega) = \frac{1}{2\pi} \intty \mean{ \bs{\mu}(0) \cdot \bs{\mu}(t) } e^{-i \omega t} dt. \label{litreview-spectral-function1}
\end{gather}

Итак, спектральная функция является Фурье-образом автокорреляционной функции оператора дипольного момента поглощающих молекул \cite{gordon1968}. При работе в единицах СИ в выражении \eqref{litreview-spectral-function} следует добавить переводной множитель
\begin{gather}
    J(\omega) = \frac{1}{4\pi \varepsilon_0} \frac{1}{2\pi} \intty \mean{ \bs{\mu}(0) \cdot \bs{\mu}(t) } e^{-i \omega t} dt. \label{litreview-spectral-function} 
\end{gather}
Следует подчеркнуть, что никаких приближений, касающихся характера движения диполей, в этом рассмотрении сделано не было. Движение системы полностью обусловлено уравнениями движения, определяемыми гамильтонианом системы $\hat{H}_0$. \par
Будем использовать следующее выражение для нормированного коэффициента поглощения через спектральную функцию 
\begin{gather}
    \alpha(\nu) = \frac{(2 \pi)^3 N_L^2}{3 \hbar c} \nu \lsq 1 - \exp \lb -\frac{h c \nu}{k T} \rb \rsq V J(\nu), \label{part1-absorption-coefficient-definition1}  
\end{gather}
%
где через $N_L$ обозначена постоянная Лошмидта, для частоты в обратных сантиметрах применяем обозначение $\nu$. 

\section{Теория временных функций корреляции и спектральных моментов} \label{section:correlation_functions}

Теория корреляционных функций получила широкое развитие для описания неравновесных систем \cite{zwanzig1965}, однако ее применение к равновесным системам также является эффективным. В системах, находящихся в термодинамическом равновесии, макроскопические параметры не претерпевают эволюции во времени, таким образом для них не имеет смысла вводить какой-то точки отсчета времени. Часто рассматривают условные вероятности, такие как $P(B, t_2 \vert A, t_1) dB$ -- вероятность того, что динамическая переменная $B$ примет значение в диапазоне $(B, \dots, B + dB)$ в момент времени $t_2$ при условии, что другая динамическая переменная имела значение $A$ в момент времени $t_1$ \cite{nitzan2006}. Также можно рассмотреть совместную вероятность $P(B, t_2; A, t_1) dB dA$ -- вероятность того, что переменная $A$ примет значение в диапазоне $(A, \dots, A+dA)$ в момент времени $t_1$ и переменная $B$ примет значение в диапазоне $(B, \dots B+dB)$ в момент времени $t_2$. Эти две вероятности связаны соотношением (формулой полной вероятности)
\begin{gather}
    P(B, t_2; A, t_1) = P(B, t_2 \vert A, t_1) P(A, t_1), 
\end{gather}
где $P(A, t_1) dA$ --вероятность того, что переменная $A$ примет значение в диапазоне $(A, \dots, A+dA)$ в момент времени $t_1$. В стационарных системах последняя вероятность, очевидно, не зависит от времени $P(A, t_1) = P(A)$; условная и совместная вероятности зависят только от разности времён
\begin{gather}
    P(B, t_2; A, t_1) = P(B, t_2 - t_1; A, 0), \quad P(B, t_2 \vert A, t_1) = P(B, t_2 - t_1 \vert A, 0),
\end{gather}
%
где $t = 0$ было положено произвольным образом.\par
Временную корреляционную функцию двух динамических переменных $A$ и $B$ определяют, как интеграл следующего вида 
\begin{gather}
    C_{AB}(t_1, t_2) = \mean{ A(t_1) B(t_2) } = \iint dA dB AB P(B, t_2; A, t_1). \label{part1-correlation-function-definition}
\end{gather}

В стационарных системах функция корреляции есть функция разности времен
\begin{gather}
    \mean{ A(t_1) B(t_2) } = \mean{ A(0) B(t) } = \mean{ A(-t) B(0) }, \quad t = t_2 - t_1.
\end{gather}

С точки зрения классической механики динамические переменные $A$, $B$ являются функциями координат $\mf{r}(t)$ и импульсов $\mf{p}(t)$ всех частиц системы
\begin{gather}
    A(t) = A \lb \mf{r}(t), \mf{p}(t) \rb, \quad B(t) = B \lb \mf{r}(t), \mf{p}(t) \rb.
\end{gather}

Фазовая траектория $\mf{r}(t)$, $\mf{p}(t)$ однозначно определена начальными условиями $\mf{r}(0)$, $\mf{p}(0)$. Таким образом, совместная вероятность в \eqref{part1-correlation-function-definition} определяется функцией распределения $f(\mf{r}, \mf{p})$ начальных условий для фазовых траекторий
\begin{gather}
    C_{AB}(t_1, t_2) = \int d\mf{r} \, d\mf{p} \, f \lb \mf{r}, \mf{p} \rb A\lb t_1; \mf{r}, \mf{p}, t = 0 \rb B \lb t_2; \mf{r}, \mf{p}, t = 0 \rb;
\end{gather}
%
обозначение $A(t_1; \mf{r}, \mf{p}, t = 0)$ означает, что динамическая переменная $A$ в момент времени $t_1$ вычисляется как функция координат и импульсов $A(\mf{r}(t_1), \mf{p}(t_1))$, вычисленных в момент времени $t_1$. \par
    При $t \rightarrow 0$ корреляционная функция $C_{AB}(t)$ становится средним значением произведения динамических переменных $A$ и $B$
\begin{gather}
    C_{AB}(0) = \mean{AB}.
\end{gather}
%
В другом пределе $t \rightarrow \infty$ можно предположить, что корреляция между переменными исчезает, то есть
\begin{gather}
    \lim_{t \rightarrow \infty} C_{AB}(t) = \mean{A} \mean{B}.
\end{gather}

Часто обьем системы, появляющийся в формуле \eqref{part1-absorption-coefficient-definition1}, рассматривают как часть спектральной функции $J(\omega)$ или рассматривают их произведение совместно. Поэтому в соответствии с обозначениями, используемыми в \cite{frommhold}, определим автокорреляционную функцию дипольного момента $C(t)$ как
\begin{gather}
    C(t) = \frac{1}{4 \pi \varepsilon_0} \frac{V}{2\pi} \mean{\bs{\mu}(0) \cdot \bs{\mu}(t)}. \label{litreview-autocorrelation-function-definition}
\end{gather}

Спектральная функция \eqref{litreview-spectral-function} принимает действительные значения, следовательно
\begin{gather}
    V J^+(\omega) = \intty C(t)^+ e^{i \omega t} dt = \intty C(-t)^+ e^{-i\omega t} dt = VJ(\omega), \label{litreview-time-inversion}
\end{gather}
%
где была сделана замена переменной $t \rightarrow -t$, а индекс $+$ обозначает комплексное сопряжение. Сранивая \eqref{litreview-time-inversion} с \eqref{litreview-spectral-function}, получаем
\begin{gather}
    C(t) = C(-t)^+.
\end{gather}

Если мы разложим корреляционную функцию на действительную и мнимую часть, то получим следующие соотношения 
\begin{gather}
    \text{Re} \, C(t) = \text{Re} \, C(-t), \quad \text{Im} \, C(t) = -\text{Im} \, C(t).
\end{gather}

То есть, действительная часть корреляционной функции является четной функции времени, а мнимая часть -- нечетной. Так как классические корреляционные функции являются действительными функциями, то они должны быть четными функциями времени. Квантово-механические же корреляционные функции, как правило, являются комплексными функциями, мнимая часть является исключительно квантово-механическим вкладом, отсутствующим при классическом рассмотрении.


Как уже отмечалось, корреляционные функции в равновесных системах зависят только от разности времени $t_1 - t_2$. Следовательно,
\begin{gather}
    0 = \frac{d}{ds} \mean{ A(t+s) B(s) } = \mean{ \dot{A}(t+s) B(s) } + \mean{ A(t+s) \dot{B}(s) } = \mean{ \dot{A}(t) B(0) } + \mean{ A(t) \dot{B}(0) }.
\end{gather}

Получаем следующее соотношение 
\begin{gather}
    \mean{\dot{A}(t)B(0)} = -\mean{A(t)\dot{B}(0)}, 
\end{gather}
которое для автокорреляционных функций переходит в 
\begin{gather}
    \mean{A \dot{A}} = 0. \label{litreview-temp1}
\end{gather}

Разложим автокорреляционную функцию в ряд по степеням времени $t$
\begin{gather}
    \mean{A(0) A(t)} = \Bigg\langle A(0) \lsq A(0) + t \dot{A}(0) + \frac{t^2}{2!} \ddot{A}(0) + \dots \rsq \Bigg\rangle = \mean{A(0) A(0)} + t \mean{A(0) \dot{A}(0)} + \frac{t^2}{2!} \mean{A(0) \ddot{A}(0)} + \dots \label{litreview-correlation-function-series}
\end{gather}

Отметим, что коэффициенты перед степенями $t^n$ не требуют знания динамики $A(t)$, а являются средними по ансамблю, т.к. временные производные могут быть записаны через скобку Пуассона
\begin{gather}
    \frac{dA}{dt} = \lsq A, H \rsq = \sum_k \lsq \frac{\partial A}{\partial x_k} \frac{\partial H}{\partial p_k} - \frac{\partial A}{\partial p_k} \frac{\partial H}{\partial x_k} \rsq. \label{litreview-poisson-bracket}
\end{gather}

Соотношение \eqref{litreview-temp1} соотносится с тем, что корреляционная функция является четной функцией. Из четности функции следует, что все коэффициенты перед нечетными степенями $t$ в \eqref{litreview-correlation-function-series} обращаются в нуль. В квантово-механической корреляционной функции нечетные степени не исчезают, и, более того, именно за счет них корреляционная функция обладает мнимой частью. \par
Коэффициенты в ряду по степеням времени \eqref{litreview-correlation-function-series} имеют физический смысл моментов соответствующего частотного спектра. Разрешим соотношение \eqref{litreview-spectral-function} относительно автокорреляционной функции дипольного момента и разложим комплексную экспоненту в подынтегральной функции в ряд
\begin{gather}
    C(t) = 2\pi V \intty J(\omega) e^{i \omega t} d \omega = 2 \pi V \intty J(\nu) e^{i 2\pi \nu c t} d\nu = 2 \pi \sum_{n = 0}^\infty \frac{\lb 2\pi i c t\rb^n}{n!} \intty \nu^n V J(\nu) d\nu. 
\end{gather}
%
Величины
\begin{gather}
    M_n = \intty \nu^n V J(\nu) d\nu \label{litreview-spectral-function-moments}
\end{gather}
%
называют $n$-ыми спектральными моментами. Теоретически знание всех спектральных моментов эквивалентно знанию спектральной функции (так называемая проблема моментов), однако на практике моменты выше второго находят редко. \par
Итак, спектральные моменты являются моментами спектральной функции и пропорциональны производным автокорреляционной функции дипольного момента в точке $t = 0$:
\begin{gather}
    \frac{1}{2 \pi} \frac{d^n C}{dt^n} \Bigg\vert_{t = 0} = \lb 2 \pi ic \rb^n M_n,
\end{gather}
%
что позволяет вычислить их как средние значения по фазовому пространству. Подробно рассматривать вопрос вычисления спектральных моментов не будем, приведем лишь выражения, по которым могут быть рассчитаны первые два спектральных момента:
\begin{gather}
    M_0 = 2 \pi \frac{\displaystyle \int \bs{\mu}^2 \exp \lsq -H(\mf{q}, \mf{p}) / \kb T \rsq d\mf{q} \, d\mf{p}}{\displaystyle \int \exp \lsq -H(\mf{q}, \mf{p}) / \kb T \rsq d\mf{q} \, d\mf{p}}, \label{litreview-m0-phase-space} \\
    M_2 = 2 \pi (2 \pi c)^2 \frac{\displaystyle \int \dot{\bs{\mu}}^2 \exp \lsq -H(\mf{q}, \mf{p}) / \kb T \rsq d\mf{q} \, d\mf{p}}{\displaystyle \int \exp \lsq -H(\mf{q}, \mf{p}) / \kb T \rsq d\mf{q} \, d\mf{p}}, \label{litreview-m2-phase-space}
\end{gather}
%
где производную по времени дипольного момента $\dot{\bs{\mu}}$ в выражении для второго спектрального момента следует преобразовать в скобку Пуассона \eqref{litreview-poisson-bracket}. Затем полученные значения интегралов следует привести к размерности см$^{-1} \cdot$Амага$^{-2}$ и  см$^{-3} \cdot$Амага$^{-2}$, соответственно. \par
Спектральные моменты, полученные по формулам \eqref{litreview-m0-phase-space}, \eqref{litreview-m2-phase-space}, могут применяться для контроля сходимости траекторных расчетов. Моменты спектральной функции $J(\omega)$, полученной в результате траекторного расчета, в пределе должны совпасть с моментами по фазовому пространству.

\newpage
\printbibliography

\end{document}
